\chapterbackmatter{Supplementary Exercises}

\begin{enumerate}
  \SupplementaryExercise{1}{A Free-Field Operator-Product Expansion}
  {Adapted from Hong Liu, MIT 8.324 Assignment 6, Problem 3 (2010).}
  {mit-8324-f10-ps6-ope-rg-v2ch20-p01}
  \begin{enumerate}
    \item[(a)]
    In four-dimensional Euclidean free scalar theory with
    \(\langle\phi(x)\phi(0)\rangle=1/(4\pi^2x^2)\), use Wick's theorem to
    find the OPE of \(\phi(x)\phi(0)\) through operators of dimension two,
    including the first derivative descendant.
    \item[(b)]
    Let \(O(x)=:\!\phi^2(x)\!:\) in the same free theory.  Determine the
    coefficients of \(\mathbf1\), \(:\!\phi^2(0)\!:\), and
    \(\partial_\mu:\!\phi^2(0)\!:\) in \(O(x)O(0)\) through the first
    descendant order.
    \item[(c)]
    For \(O(x)=:\!\phi^4(x)\!:\), use cross-contractions to list all scalar
    normal-ordered operators of canonical dimension at most four that
    occur in \(O(x)O(0)\), and determine the leading coefficient of each.
  \end{enumerate}

  \SupplementaryExercise{2}{An \(O(N)\) OPE and the Wilson--Fisher Exponent}
  {Adapted from John McGreevy, Physics 215C, Homework 5, Problem 2 (2025).}
  {mcgreevy-215c-s25-hw5-q2-v2ch20-p02}
  \begin{enumerate}
    \item[(a)]
    For \(N\) free real scalars with
    \(\langle\phi_i(x)\phi_j(0)\rangle=\delta_{ij}G(x)\), let
    \(E=:\!\phi_i\phi_i\!:\).  Use Wick's theorem to determine the
    coefficients of \(\mathbf1\) and \(E\) in \(E(x)E(0)\).
    \item[(b)]
    Perturb the theory in \(d=4-\epsilon\) by
    \(\lambda(\phi_i\phi_i)^2/4!\) and define
    \(u=\lambda/(8\pi^2)\).  Use the one-loop OPE contractions to obtain
    \[
      \beta(u)=-\epsilon u+\frac{N+8}{6}u^2+O(u^3)
    \]
    and find the nontrivial fixed point.
    \item[(c)]
    The thermal perturbation \(tE/2\) has eigenvalue
    \(y_t=2-(N+2)u/6+O(u^2)\).  Evaluate it at the fixed point and derive
    the correlation-length exponent \(\nu=1/y_t\) through
    \(O(\epsilon)\).
    \item[(d)]
    Check the \(N=1\) result and explain how the free OPE data control both
    operator mixing and the leading interacting critical exponent.
  \end{enumerate}

  \SupplementaryExercise{3}{Scheme Changes and Dimensional Transmutation}
  {Adapted from Hong Liu, MIT 8.324 Assignment 6, Problem 4 (2010).}
  {mit-8324-f10-ps6-ope-rg-v2ch20-p03}
  \begin{enumerate}
    \item[(a)] Treat the beta function as a flow on coupling space.  Under
      an invertible redefinition \(\bar\lambda(\lambda)\), derive its exact
      transformation law and identify its tensor character.
    \item[(b)] Let
      \[
        \beta(\lambda)=b_2\lambda^2+b_3\lambda^3+b_4\lambda^4+\cdots,
        \qquad
        \bar\lambda=\lambda+a_2\lambda^2+a_3\lambda^3+\cdots .
      \]
      Show that \(\bar b_2=b_2\) and \(\bar b_3=b_3\), exhibit the freedom
      to choose \(\bar b_4\), prove that fixed points map to fixed points,
      and compare \(\beta'\) and \(\bar\beta'\) at a fixed point.
    \item[(c)] For
      \[
        \mu\frac{dg}{d\mu}=-bg^2-cg^3-dg^4-\cdots ,
      \]
      derive
      \[
        \log\frac{\mu}{\Lambda}
          =\frac1{bg(\mu)}+\frac{c}{b^2}\log[bg(\mu)]+O(g),
      \]
      and prove that \(\Lambda\) is RG invariant.
    \item[(d)] In a renormalizable theory whose only parameter is one
      dimensionless running coupling, show generally that
      \(g(\mu)=f(\log(\mu/\Lambda))\).  Explain dimensional transmutation.
  \end{enumerate}

  \SupplementaryExercise{4}{Current Correlators, Dispersion Relations, and Finite-Energy Sum Rules}
  {Inspired by A. V. Radyushkin, Introduction to QCD Sum Rule Approach
  (2001), Sections 3--4.}
  {radyushkin-qcd-sum-rules-2001-v2ch20-p04}
  \begin{enumerate}
    \item[(a)]
    For a conserved Hermitian current \(J^\mu\), show that its vacuum
    correlator has the form
    \(\Pi^{\mu\nu}(q)=(q^\mu q^\nu-q^2\eta^{\mu\nu})\Pi(q^2)\), up to
    contact terms.  State the positivity property of its spectral density.
    \item[(b)]
    Starting from
    \(\Pi(Q^2)=\int_0^\infty ds\,\rho(s)/(s+Q^2)\), expand for large
    Euclidean \(Q^2\).  Relate the first three inverse powers of \(Q^2\) to
    moments of \(\rho\), and state when this termwise expansion is valid.
    \item[(c)]
    Let \(\Pi(s)\) be analytic except for a cut on the positive real axis.
    Apply Cauchy's theorem to \(s^n\Pi(s)\) on a circle \(|s|=s_0\) and
    derive a finite-energy sum rule for
    \(\int_0^{s_0}ds\,s^n\rho(s)\).
  \end{enumerate}

  \SupplementaryExercise{5}{Borel Sum Rules, Spectral Models, and Condensate Power Counting}
  {Inspired by A. V. Radyushkin, Introduction to QCD Sum Rule Approach
  (2001), Sections 2--4.}
  {radyushkin-qcd-sum-rules-2001-v2ch20-p05}
  \begin{enumerate}
    \item[(a)]
    Apply the QCD-sum-rule Borel operator to
    \(1/(s+Q^2)\) and show that it produces an exponential weight.
    Explain why this suppresses a high-energy continuum relative to a
    narrow low-mass resonance and removes polynomial subtraction terms.
    \item[(b)]
    Take
    \(\rho(s)=f_R^2\delta(s-m_R^2)+\theta(s-s_0)\rho_{\rm pert}(s)\).
    Write its Borel-weighted sum rule and show how differentiating with
    respect to \(1/M^2\) can isolate \(m_R^2\) as a ratio of moments.
    \item[(c)]
    A correlator of two dimension-three currents has scalar invariant
    \(\Pi(Q^2)\) of dimension zero.  Determine the large-\(Q^2\) power
    multiplying a vacuum condensate of dimension \(d\), and compare the
    dimension-four gluon and dimension-six four-quark contributions.
  \end{enumerate}

  \SupplementaryExercise{6}{Deep-Inelastic Scattering from a Photon}
  {Adapted from Iain Stewart, MIT 8.325 QFT III, Homework 7, Problem 3
  (2007), which assigns Peskin--Schroeder Problem 18.4(a)--(c).}
  {mit-8325-s07-dis-v2ch20-p06}
  \begin{enumerate}
    \item[(a)] Use the QED splitting kernels
      \[
      P_{e\leftarrow\gamma}=z^2+(1-z)^2,\quad
      P_{\gamma\leftarrow e}=\frac{1+(1-z)^2}{z},
      \]
      together with the standard plus-prescribed
      \(P_{e\leftarrow e}\) and \(P_{\gamma\leftarrow\gamma}\), to write
      the coupled evolution equations for \(f_q^\gamma,f_{\bar q}^\gamma\),
      and \(f_\gamma^\gamma\).  Starting from a pointlike photon at \(Q_0\),
      solve through first order in \(\alpha\).
    \item[(b)] Define the forward current product between photon states and
      decompose it into two transverse structure functions.  Give its
      leading-twist OPE moment expansion and obtain the one-loop mixing
      coefficients among quark and photon twist-two operators.
    \item[(c)] For \(n=2\), two flavors \(u,d\), and three colors, construct
      the anomalous-dimension matrix in the basis
      \((O_u,O_d,O_\gamma)\).  Explain how this matrix governs the second
      moment of the photon structure function.
  \end{enumerate}

  \SupplementaryExercise{7}{Parton Sum Rules and Plus Distributions}
  {Inspired by the Particle Data Group, Structure Functions review
  (2025), Sections 18.1--18.2.}
  {pdg-structure-functions-2025-v2ch20-p07}
  \begin{enumerate}
    \item[(a)]
    Define \(q_v(x)=q(x)-\bar q(x)\).  State and justify the proton sum
    rules for the integrals of \(u_v,d_v\), and all other light-flavor
    valence distributions.  Identify the conserved local operators behind
    these relations.
    \item[(b)]
    Write the quark, antiquark, and gluon momentum sum rule at a
    factorization scale \(\mu\).  Explain why QCD evolution changes the
    individual terms but cannot change their sum.
    \item[(c)]
    Using
    \(\int_0^1dx\,f(x)/(1-x)_+
      =\int_0^1dx\,[f(x)-f(1)]/(1-x)\),
    evaluate the zeroth and first moments of
    \((1+x^2)/(1-x)_+\).  Explain the role of an accompanying
    \(\delta(1-x)\) term.
  \end{enumerate}

  \SupplementaryExercise{8}{Mellin Evolution, Conservation Constraints, and Higher Twist}
  {Inspired by the Particle Data Group, Structure Functions review
  (2025), Section 18.2.}
  {pdg-structure-functions-2025-v2ch20-p08}
  \begin{enumerate}
    \item[(a)]
    For
    \((P\otimes f)(x)=\int_x^1dy\,P(x/y)f(y)/y\), prove that its
    \(n\)th Mellin moment equals the product of the \(n\)th moments of
    \(P\) and \(f\).  State the resulting simplification of DGLAP evolution.
    \item[(b)]
    Derive the Mellin-moment constraints on nonsinglet and singlet splitting
    kernels that preserve valence number and total momentum.  Explain why
    these constraints require virtual \(\delta(1-x)\) contributions.
    \item[(c)]
    Suppose a structure-function moment is fitted by
    \(M(Q^2)=M_{\rm LT}(Q^2)+H/Q^2+K/Q^4\).
    Assign the twists of the two power corrections, and explain why they
    become most visible at modest \(Q^2\) or near \(x=1\).
  \end{enumerate}

  \SupplementaryExercise{9}{Borel Singularities and Renormalon Ambiguities}
  {Inspired by M. Beneke, Renormalons (1999), Sections 1--2.}
  {beneke-renormalons-1999-v2ch20-p09}
  \begin{enumerate}
    \item[(a)]
    For coefficients \(f_n=K\,n!\,a^{-n}\), compute the Borel transform
    \(B(z)=\sum_nf_nz^n/n!\).  Locate its singularity and determine whether
    the positive-coupling Borel integral is ambiguous for positive and
    negative \(a\).
    \item[(b)]
    Compare series with large-order behavior
    \(f_n\sim(-1)^n n!A^n\) and \(f_n\sim n!A^n\), with \(A>0\).
    Locate their leading Borel singularities and explain which series is
    Borel summable without a contour prescription.
    \item[(c)]
    If \(B(z)=R/(z_0-z)\) with \(z_0>0\), evaluate the difference between
    contours passing above and below the pole in the Borel integral.
    Show its nonperturbative dependence on the coupling.
  \end{enumerate}

  \SupplementaryExercise{10}{Renormalons, Optimal Truncation, and OPE Cancellation}
  {Inspired by M. Beneke, Renormalons (1999), Sections 3 and 6.}
  {beneke-renormalons-1999-v2ch20-p10}
  \begin{enumerate}
    \item[(a)]
    In an asymptotically free theory,
    \(\Lambda/Q\sim\exp[-1/(2\beta_0\alpha(Q))]\).
    Show how a Borel singularity at \(u=p/2\) produces an ambiguity of
    order \((\Lambda/Q)^p\), and identify the required OPE power correction.
    \item[(b)]
    For \(f_n\sim n!A^n\), estimate the order \(n_*\) of the smallest term
    in \(\sum_nf_n\alpha^n\) at small positive \(\alpha\).  Estimate the
    irreducible remainder and relate it to the nearest Borel singularity.
    \item[(c)]
    Explain why a physical quantity cannot retain the contour ambiguity of
    a perturbative Wilson coefficient.  Show schematically how a change in
    the prescription for a dimension-\(d\) coefficient is cancelled by
    the corresponding prescription dependence of an OPE matrix element.
  \end{enumerate}

\end{enumerate}
